%!TEX root = ../../main.tex
\section{M-RQ2에 대한 답}
\label{sec:value-answer}

전략적 가치 개선은 경기 승패로 학습해 동결한 평가기의 교전 구간 전후 변화의 부호로 정의된다(식 \eqref{eq:svi}). 이 측정은 다섯 가지로 특성화된다. 평가기는 실제 승패를 순위화하며 그 품질은 시간대에 따라 다르다(C1--C3). 변화의 방향·평균·규모는 서로 다른 객체이다(C4). 교전 구간의 절대 변화는 짝지은 비교전 구간보다 크다(C5). 결정된 킬 축과의 부호 일치는 0.9 수준이고 다른 물질 축과의 대응은 그보다 약하며, 교환가치 라벨과는 \meFlip{}의 사례에서 부호가 갈린다(C6, C7, C31). 지평을 바꿔도 같은 결과 시점을 공유하는 행에서 라벨은 유지되고, 동료 평가기와의 부호 일치는 \peerSignAgree{}이나 작은 변화량에서는 예측 비교의 부호가 바뀐다(C8, C29). 구간 변화의 산술 분해에서 관측 갱신 성분의 절대 규모가 시간 경과 성분보다 크다(C25). 모든 값은 동결된 fit85 평가기와 결과 시점 규칙에 조건부이며, 각 결과가 무엇을 재는지의 조건은 해당 절에, 공통 해석 범위는 제 \ref{sec:disc-scope} 절에 둔다.
